% ======================================================================
% Estilo común de los manuales del Sistema de Biblioteca UMAG.
% Lo cargan manual-usuario.tex y manual-admin.tex con % ======================================================================
% Estilo común de los manuales del Sistema de Biblioteca UMAG.
% Lo cargan manual-usuario.tex y manual-admin.tex con % ======================================================================
% Estilo común de los manuales del Sistema de Biblioteca UMAG.
% Lo cargan manual-usuario.tex y manual-admin.tex con % ======================================================================
% Estilo común de los manuales del Sistema de Biblioteca UMAG.
% Lo cargan manual-usuario.tex y manual-admin.tex con \input{estilo-manual},
% después de definir \nombremanual y \versionmanual.
% ======================================================================

% ---------- Codificación e idioma ----------
\usepackage[utf8]{inputenc}
\usepackage[T1]{fontenc}
\usepackage[spanish,es-noindentfirst]{babel}

% ---------- Diseño de página ----------
\usepackage[a4paper,margin=2.5cm,headheight=15pt]{geometry}
\usepackage{xcolor}
\usepackage{graphicx}
\usepackage{fancyhdr}
\usepackage{titlesec}
\usepackage{enumitem}
\usepackage[most]{tcolorbox}
\usepackage{booktabs}
\usepackage{tabularx}
\usepackage{caption}
\usepackage{needspace}
\usepackage{etoolbox}
\usepackage{amssymb}
\usepackage{hyperref}

% ---------- Paleta institucional (la misma de la app: docs/estilo-visual.md) ----------
% Morados del degradado de la barra superior y dorado de la pestaña activa.
\definecolor{umagpurple}{HTML}{3B28A3}
\definecolor{umagpurpledark}{HTML}{2D1B69}
\definecolor{umaggold}{HTML}{B08D57}
\definecolor{umaggray}{HTML}{4B5563}
\definecolor{cajafondo}{HTML}{F8FAFC}

\hypersetup{
  colorlinks=true,
  linkcolor=umagpurple,
  urlcolor=umagpurple,
  citecolor=umagpurple,
  pdftitle={\nombremanual{} - Sistema de Biblioteca UMAG},
  pdfauthor={Biblioteca UMAG},
}

% ---------- Encabezados / pies de página ----------
\pagestyle{fancy}
\fancyhf{}
\renewcommand{\headrulewidth}{0.4pt}
\renewcommand{\footrulewidth}{0pt}
% Encabezado de capítulo sin el prefijo "Capítulo N." (ya se ve grande en la
% portadilla del capítulo) para que quepa junto al nombre del manual sin superponerse.
\renewcommand{\chaptermark}[1]{\markboth{#1}{}}
\fancyhead[L]{\small\textcolor{umaggray}{Biblioteca UMAG · \nombremanual}}
\fancyhead[R]{\small\textcolor{umaggray}{\leftmark}}
\fancyfoot[C]{\small\textcolor{umaggray}{\thepage}}

\fancypagestyle{plain}{%
  \fancyhf{}
  \fancyfoot[C]{\small\textcolor{umaggray}{\thepage}}
  \renewcommand{\headrulewidth}{0pt}
}

% ---------- Estilo de títulos ----------
\titleformat{\chapter}[display]
  {\normalfont\bfseries\color{umagpurple}}
  {\Large\filright\textcolor{umaggold}{CAPÍTULO \thechapter}}
  {6pt}
  {\Huge\filright}
\titlespacing*{\chapter}{0pt}{0pt}{28pt}

\titleformat{\section}
  {\normalfont\Large\bfseries\color{umagpurpledark}}{\thesection}{0.6em}{}
\titleformat{\subsection}
  {\normalfont\large\bfseries\color{umagpurple}}{\thesubsection}{0.6em}{}
% Un título de sección no queda solo al pie de la página: si no caben él,
% su etiqueta de permiso y unas líneas de texto, pasa a la página siguiente.
\pretocmd{\section}{\needspace{7\baselineskip}}{}{}

% ---------- Cajas reutilizables ----------
\newtcolorbox{nota}{
  colback=umaggold!8, colframe=umaggold, boxrule=0.6pt, arc=2pt,
  left=8pt, right=8pt, top=6pt, bottom=6pt,
  fonttitle=\bfseries, title={Tip},
}
\newtcolorbox{importante}{
  colback=umagpurple!6, colframe=umagpurple, boxrule=0.6pt, arc=2pt,
  left=8pt, right=8pt, top=6pt, bottom=6pt,
  fonttitle=\bfseries, title={Importante},
}
% Etiqueta compacta (no una caja que envuelve contenido) para marcar una
% sección que solo ven los roles con cierto permiso, sin arriesgar que una
% tcolorbox con una captura de pantalla adentro desborde el alto de la página.
% \requierepermiso{nombre del permiso, tal como aparece en Personal y Roles}
\newtcbox{\etiquetapermiso}{on line, colback=umagpurple, colframe=umagpurple,
  coltext=white, boxrule=0pt, arc=0pt, left=4pt, right=4pt, top=2pt, bottom=2pt,
  fontupper=\bfseries\small}
\newcommand{\requierepermiso}[1]{%
  \needspace{3\baselineskip}%
  \par\noindent\etiquetapermiso{REQUIERE PERMISO: #1}\par\vspace{6pt}%
}

% ---------- Comando para insertar capturas de pantalla ----------
% \captura{archivo}{texto de pie de figura}
\newcommand{\captura}[2]{%
  \begin{center}
  \fbox{\includegraphics[width=0.92\textwidth,height=0.78\textheight,keepaspectratio]{imagenes/#1}}\\[5pt]
  {\small\itshape\color{umaggray} #2}
  \end{center}
  \vspace{6pt}
}

\setlist[itemize]{leftmargin=1.4em}
\setlist[enumerate]{leftmargin=1.6em}

% ---------- Portada ----------
% \portadamanual{subtítulo}{contenido del recuadro de acceso}
\newcommand{\portadamanual}[2]{%
\begin{titlepage}
\centering
\vspace*{1.5cm}

{\includegraphics[width=3cm]{imagenes/logo-umag.png}}

\vspace{0.8cm}

{\color{umaggray}\large UNIVERSIDAD DE MAGALLANES}

\vspace{2.2cm}

{\color{umagpurple}\Huge\bfseries Sistema de Biblioteca}\\[6pt]
{\color{umagpurpledark}\Huge\bfseries UMAG}

\vspace{1cm}
{\color{umaggold}\rule{8cm}{1.2pt}}
\vspace{1cm}

{\color{umaggray}\LARGE \nombremanual}

\vspace{0.4cm}
{\color{umaggray}\large #1}

\vfill

\begin{tcolorbox}[colback=cajafondo,colframe=umagpurple!40,width=0.8\textwidth,arc=3pt,boxrule=0.6pt]
\centering
#2
\end{tcolorbox}

\vspace{1cm}
{\color{umaggray} Versión del documento: \versionmanual}

\vspace{1cm}
\end{titlepage}
}
,
% después de definir \nombremanual y \versionmanual.
% ======================================================================

% ---------- Codificación e idioma ----------
\usepackage[utf8]{inputenc}
\usepackage[T1]{fontenc}
\usepackage[spanish,es-noindentfirst]{babel}

% ---------- Diseño de página ----------
\usepackage[a4paper,margin=2.5cm,headheight=15pt]{geometry}
\usepackage{xcolor}
\usepackage{graphicx}
\usepackage{fancyhdr}
\usepackage{titlesec}
\usepackage{enumitem}
\usepackage[most]{tcolorbox}
\usepackage{booktabs}
\usepackage{tabularx}
\usepackage{caption}
\usepackage{needspace}
\usepackage{etoolbox}
\usepackage{amssymb}
\usepackage{hyperref}

% ---------- Paleta institucional (la misma de la app: docs/estilo-visual.md) ----------
% Morados del degradado de la barra superior y dorado de la pestaña activa.
\definecolor{umagpurple}{HTML}{3B28A3}
\definecolor{umagpurpledark}{HTML}{2D1B69}
\definecolor{umaggold}{HTML}{B08D57}
\definecolor{umaggray}{HTML}{4B5563}
\definecolor{cajafondo}{HTML}{F8FAFC}

\hypersetup{
  colorlinks=true,
  linkcolor=umagpurple,
  urlcolor=umagpurple,
  citecolor=umagpurple,
  pdftitle={\nombremanual{} - Sistema de Biblioteca UMAG},
  pdfauthor={Biblioteca UMAG},
}

% ---------- Encabezados / pies de página ----------
\pagestyle{fancy}
\fancyhf{}
\renewcommand{\headrulewidth}{0.4pt}
\renewcommand{\footrulewidth}{0pt}
% Encabezado de capítulo sin el prefijo "Capítulo N." (ya se ve grande en la
% portadilla del capítulo) para que quepa junto al nombre del manual sin superponerse.
\renewcommand{\chaptermark}[1]{\markboth{#1}{}}
\fancyhead[L]{\small\textcolor{umaggray}{Biblioteca UMAG · \nombremanual}}
\fancyhead[R]{\small\textcolor{umaggray}{\leftmark}}
\fancyfoot[C]{\small\textcolor{umaggray}{\thepage}}

\fancypagestyle{plain}{%
  \fancyhf{}
  \fancyfoot[C]{\small\textcolor{umaggray}{\thepage}}
  \renewcommand{\headrulewidth}{0pt}
}

% ---------- Estilo de títulos ----------
\titleformat{\chapter}[display]
  {\normalfont\bfseries\color{umagpurple}}
  {\Large\filright\textcolor{umaggold}{CAPÍTULO \thechapter}}
  {6pt}
  {\Huge\filright}
\titlespacing*{\chapter}{0pt}{0pt}{28pt}

\titleformat{\section}
  {\normalfont\Large\bfseries\color{umagpurpledark}}{\thesection}{0.6em}{}
\titleformat{\subsection}
  {\normalfont\large\bfseries\color{umagpurple}}{\thesubsection}{0.6em}{}
% Un título de sección no queda solo al pie de la página: si no caben él,
% su etiqueta de permiso y unas líneas de texto, pasa a la página siguiente.
\pretocmd{\section}{\needspace{7\baselineskip}}{}{}

% ---------- Cajas reutilizables ----------
\newtcolorbox{nota}{
  colback=umaggold!8, colframe=umaggold, boxrule=0.6pt, arc=2pt,
  left=8pt, right=8pt, top=6pt, bottom=6pt,
  fonttitle=\bfseries, title={Tip},
}
\newtcolorbox{importante}{
  colback=umagpurple!6, colframe=umagpurple, boxrule=0.6pt, arc=2pt,
  left=8pt, right=8pt, top=6pt, bottom=6pt,
  fonttitle=\bfseries, title={Importante},
}
% Etiqueta compacta (no una caja que envuelve contenido) para marcar una
% sección que solo ven los roles con cierto permiso, sin arriesgar que una
% tcolorbox con una captura de pantalla adentro desborde el alto de la página.
% \requierepermiso{nombre del permiso, tal como aparece en Personal y Roles}
\newtcbox{\etiquetapermiso}{on line, colback=umagpurple, colframe=umagpurple,
  coltext=white, boxrule=0pt, arc=0pt, left=4pt, right=4pt, top=2pt, bottom=2pt,
  fontupper=\bfseries\small}
\newcommand{\requierepermiso}[1]{%
  \needspace{3\baselineskip}%
  \par\noindent\etiquetapermiso{REQUIERE PERMISO: #1}\par\vspace{6pt}%
}

% ---------- Comando para insertar capturas de pantalla ----------
% \captura{archivo}{texto de pie de figura}
\newcommand{\captura}[2]{%
  \begin{center}
  \fbox{\includegraphics[width=0.92\textwidth,height=0.78\textheight,keepaspectratio]{imagenes/#1}}\\[5pt]
  {\small\itshape\color{umaggray} #2}
  \end{center}
  \vspace{6pt}
}

\setlist[itemize]{leftmargin=1.4em}
\setlist[enumerate]{leftmargin=1.6em}

% ---------- Portada ----------
% \portadamanual{subtítulo}{contenido del recuadro de acceso}
\newcommand{\portadamanual}[2]{%
\begin{titlepage}
\centering
\vspace*{1.5cm}

{\includegraphics[width=3cm]{imagenes/logo-umag.png}}

\vspace{0.8cm}

{\color{umaggray}\large UNIVERSIDAD DE MAGALLANES}

\vspace{2.2cm}

{\color{umagpurple}\Huge\bfseries Sistema de Biblioteca}\\[6pt]
{\color{umagpurpledark}\Huge\bfseries UMAG}

\vspace{1cm}
{\color{umaggold}\rule{8cm}{1.2pt}}
\vspace{1cm}

{\color{umaggray}\LARGE \nombremanual}

\vspace{0.4cm}
{\color{umaggray}\large #1}

\vfill

\begin{tcolorbox}[colback=cajafondo,colframe=umagpurple!40,width=0.8\textwidth,arc=3pt,boxrule=0.6pt]
\centering
#2
\end{tcolorbox}

\vspace{1cm}
{\color{umaggray} Versión del documento: \versionmanual}

\vspace{1cm}
\end{titlepage}
}
,
% después de definir \nombremanual y \versionmanual.
% ======================================================================

% ---------- Codificación e idioma ----------
\usepackage[utf8]{inputenc}
\usepackage[T1]{fontenc}
\usepackage[spanish,es-noindentfirst]{babel}

% ---------- Diseño de página ----------
\usepackage[a4paper,margin=2.5cm,headheight=15pt]{geometry}
\usepackage{xcolor}
\usepackage{graphicx}
\usepackage{fancyhdr}
\usepackage{titlesec}
\usepackage{enumitem}
\usepackage[most]{tcolorbox}
\usepackage{booktabs}
\usepackage{tabularx}
\usepackage{caption}
\usepackage{needspace}
\usepackage{etoolbox}
\usepackage{amssymb}
\usepackage{hyperref}

% ---------- Paleta institucional (la misma de la app: docs/estilo-visual.md) ----------
% Morados del degradado de la barra superior y dorado de la pestaña activa.
\definecolor{umagpurple}{HTML}{3B28A3}
\definecolor{umagpurpledark}{HTML}{2D1B69}
\definecolor{umaggold}{HTML}{B08D57}
\definecolor{umaggray}{HTML}{4B5563}
\definecolor{cajafondo}{HTML}{F8FAFC}

\hypersetup{
  colorlinks=true,
  linkcolor=umagpurple,
  urlcolor=umagpurple,
  citecolor=umagpurple,
  pdftitle={\nombremanual{} - Sistema de Biblioteca UMAG},
  pdfauthor={Biblioteca UMAG},
}

% ---------- Encabezados / pies de página ----------
\pagestyle{fancy}
\fancyhf{}
\renewcommand{\headrulewidth}{0.4pt}
\renewcommand{\footrulewidth}{0pt}
% Encabezado de capítulo sin el prefijo "Capítulo N." (ya se ve grande en la
% portadilla del capítulo) para que quepa junto al nombre del manual sin superponerse.
\renewcommand{\chaptermark}[1]{\markboth{#1}{}}
\fancyhead[L]{\small\textcolor{umaggray}{Biblioteca UMAG · \nombremanual}}
\fancyhead[R]{\small\textcolor{umaggray}{\leftmark}}
\fancyfoot[C]{\small\textcolor{umaggray}{\thepage}}

\fancypagestyle{plain}{%
  \fancyhf{}
  \fancyfoot[C]{\small\textcolor{umaggray}{\thepage}}
  \renewcommand{\headrulewidth}{0pt}
}

% ---------- Estilo de títulos ----------
\titleformat{\chapter}[display]
  {\normalfont\bfseries\color{umagpurple}}
  {\Large\filright\textcolor{umaggold}{CAPÍTULO \thechapter}}
  {6pt}
  {\Huge\filright}
\titlespacing*{\chapter}{0pt}{0pt}{28pt}

\titleformat{\section}
  {\normalfont\Large\bfseries\color{umagpurpledark}}{\thesection}{0.6em}{}
\titleformat{\subsection}
  {\normalfont\large\bfseries\color{umagpurple}}{\thesubsection}{0.6em}{}
% Un título de sección no queda solo al pie de la página: si no caben él,
% su etiqueta de permiso y unas líneas de texto, pasa a la página siguiente.
\pretocmd{\section}{\needspace{7\baselineskip}}{}{}

% ---------- Cajas reutilizables ----------
\newtcolorbox{nota}{
  colback=umaggold!8, colframe=umaggold, boxrule=0.6pt, arc=2pt,
  left=8pt, right=8pt, top=6pt, bottom=6pt,
  fonttitle=\bfseries, title={Tip},
}
\newtcolorbox{importante}{
  colback=umagpurple!6, colframe=umagpurple, boxrule=0.6pt, arc=2pt,
  left=8pt, right=8pt, top=6pt, bottom=6pt,
  fonttitle=\bfseries, title={Importante},
}
% Etiqueta compacta (no una caja que envuelve contenido) para marcar una
% sección que solo ven los roles con cierto permiso, sin arriesgar que una
% tcolorbox con una captura de pantalla adentro desborde el alto de la página.
% \requierepermiso{nombre del permiso, tal como aparece en Personal y Roles}
\newtcbox{\etiquetapermiso}{on line, colback=umagpurple, colframe=umagpurple,
  coltext=white, boxrule=0pt, arc=0pt, left=4pt, right=4pt, top=2pt, bottom=2pt,
  fontupper=\bfseries\small}
\newcommand{\requierepermiso}[1]{%
  \needspace{3\baselineskip}%
  \par\noindent\etiquetapermiso{REQUIERE PERMISO: #1}\par\vspace{6pt}%
}

% ---------- Comando para insertar capturas de pantalla ----------
% \captura{archivo}{texto de pie de figura}
\newcommand{\captura}[2]{%
  \begin{center}
  \fbox{\includegraphics[width=0.92\textwidth,height=0.78\textheight,keepaspectratio]{imagenes/#1}}\\[5pt]
  {\small\itshape\color{umaggray} #2}
  \end{center}
  \vspace{6pt}
}

\setlist[itemize]{leftmargin=1.4em}
\setlist[enumerate]{leftmargin=1.6em}

% ---------- Portada ----------
% \portadamanual{subtítulo}{contenido del recuadro de acceso}
\newcommand{\portadamanual}[2]{%
\begin{titlepage}
\centering
\vspace*{1.5cm}

{\includegraphics[width=3cm]{imagenes/logo-umag.png}}

\vspace{0.8cm}

{\color{umaggray}\large UNIVERSIDAD DE MAGALLANES}

\vspace{2.2cm}

{\color{umagpurple}\Huge\bfseries Sistema de Biblioteca}\\[6pt]
{\color{umagpurpledark}\Huge\bfseries UMAG}

\vspace{1cm}
{\color{umaggold}\rule{8cm}{1.2pt}}
\vspace{1cm}

{\color{umaggray}\LARGE \nombremanual}

\vspace{0.4cm}
{\color{umaggray}\large #1}

\vfill

\begin{tcolorbox}[colback=cajafondo,colframe=umagpurple!40,width=0.8\textwidth,arc=3pt,boxrule=0.6pt]
\centering
#2
\end{tcolorbox}

\vspace{1cm}
{\color{umaggray} Versión del documento: \versionmanual}

\vspace{1cm}
\end{titlepage}
}
,
% después de definir \nombremanual y \versionmanual.
% ======================================================================

% ---------- Codificación e idioma ----------
\usepackage[utf8]{inputenc}
\usepackage[T1]{fontenc}
\usepackage[spanish,es-noindentfirst]{babel}

% ---------- Diseño de página ----------
\usepackage[a4paper,margin=2.5cm,headheight=15pt]{geometry}
\usepackage{xcolor}
\usepackage{graphicx}
\usepackage{fancyhdr}
\usepackage{titlesec}
\usepackage{enumitem}
\usepackage[most]{tcolorbox}
\usepackage{booktabs}
\usepackage{tabularx}
\usepackage{caption}
\usepackage{needspace}
\usepackage{etoolbox}
\usepackage{amssymb}
\usepackage{hyperref}

% ---------- Paleta institucional (la misma de la app: docs/estilo-visual.md) ----------
% Morados del degradado de la barra superior y dorado de la pestaña activa.
\definecolor{umagpurple}{HTML}{3B28A3}
\definecolor{umagpurpledark}{HTML}{2D1B69}
\definecolor{umaggold}{HTML}{B08D57}
\definecolor{umaggray}{HTML}{4B5563}
\definecolor{cajafondo}{HTML}{F8FAFC}

\hypersetup{
  colorlinks=true,
  linkcolor=umagpurple,
  urlcolor=umagpurple,
  citecolor=umagpurple,
  pdftitle={\nombremanual{} - Sistema de Biblioteca UMAG},
  pdfauthor={Biblioteca UMAG},
}

% ---------- Encabezados / pies de página ----------
\pagestyle{fancy}
\fancyhf{}
\renewcommand{\headrulewidth}{0.4pt}
\renewcommand{\footrulewidth}{0pt}
% Encabezado de capítulo sin el prefijo "Capítulo N." (ya se ve grande en la
% portadilla del capítulo) para que quepa junto al nombre del manual sin superponerse.
\renewcommand{\chaptermark}[1]{\markboth{#1}{}}
\fancyhead[L]{\small\textcolor{umaggray}{Biblioteca UMAG · \nombremanual}}
\fancyhead[R]{\small\textcolor{umaggray}{\leftmark}}
\fancyfoot[C]{\small\textcolor{umaggray}{\thepage}}

\fancypagestyle{plain}{%
  \fancyhf{}
  \fancyfoot[C]{\small\textcolor{umaggray}{\thepage}}
  \renewcommand{\headrulewidth}{0pt}
}

% ---------- Estilo de títulos ----------
\titleformat{\chapter}[display]
  {\normalfont\bfseries\color{umagpurple}}
  {\Large\filright\textcolor{umaggold}{CAPÍTULO \thechapter}}
  {6pt}
  {\Huge\filright}
\titlespacing*{\chapter}{0pt}{0pt}{28pt}

\titleformat{\section}
  {\normalfont\Large\bfseries\color{umagpurpledark}}{\thesection}{0.6em}{}
\titleformat{\subsection}
  {\normalfont\large\bfseries\color{umagpurple}}{\thesubsection}{0.6em}{}
% Un título de sección no queda solo al pie de la página: si no caben él,
% su etiqueta de permiso y unas líneas de texto, pasa a la página siguiente.
\pretocmd{\section}{\needspace{7\baselineskip}}{}{}

% ---------- Cajas reutilizables ----------
\newtcolorbox{nota}{
  colback=umaggold!8, colframe=umaggold, boxrule=0.6pt, arc=2pt,
  left=8pt, right=8pt, top=6pt, bottom=6pt,
  fonttitle=\bfseries, title={Tip},
}
\newtcolorbox{importante}{
  colback=umagpurple!6, colframe=umagpurple, boxrule=0.6pt, arc=2pt,
  left=8pt, right=8pt, top=6pt, bottom=6pt,
  fonttitle=\bfseries, title={Importante},
}
% Etiqueta compacta (no una caja que envuelve contenido) para marcar una
% sección que solo ven los roles con cierto permiso, sin arriesgar que una
% tcolorbox con una captura de pantalla adentro desborde el alto de la página.
% \requierepermiso{nombre del permiso, tal como aparece en Personal y Roles}
\newtcbox{\etiquetapermiso}{on line, colback=umagpurple, colframe=umagpurple,
  coltext=white, boxrule=0pt, arc=0pt, left=4pt, right=4pt, top=2pt, bottom=2pt,
  fontupper=\bfseries\small}
\newcommand{\requierepermiso}[1]{%
  \needspace{3\baselineskip}%
  \par\noindent\etiquetapermiso{REQUIERE PERMISO: #1}\par\vspace{6pt}%
}

% ---------- Comando para insertar capturas de pantalla ----------
% \captura{archivo}{texto de pie de figura}
\newcommand{\captura}[2]{%
  \begin{center}
  \fbox{\includegraphics[width=0.92\textwidth,height=0.78\textheight,keepaspectratio]{imagenes/#1}}\\[5pt]
  {\small\itshape\color{umaggray} #2}
  \end{center}
  \vspace{6pt}
}

\setlist[itemize]{leftmargin=1.4em}
\setlist[enumerate]{leftmargin=1.6em}

% ---------- Portada ----------
% \portadamanual{subtítulo}{contenido del recuadro de acceso}
\newcommand{\portadamanual}[2]{%
\begin{titlepage}
\centering
\vspace*{1.5cm}

{\includegraphics[width=3cm]{imagenes/logo-umag.png}}

\vspace{0.8cm}

{\color{umaggray}\large UNIVERSIDAD DE MAGALLANES}

\vspace{2.2cm}

{\color{umagpurple}\Huge\bfseries Sistema de Biblioteca}\\[6pt]
{\color{umagpurpledark}\Huge\bfseries UMAG}

\vspace{1cm}
{\color{umaggold}\rule{8cm}{1.2pt}}
\vspace{1cm}

{\color{umaggray}\LARGE \nombremanual}

\vspace{0.4cm}
{\color{umaggray}\large #1}

\vfill

\begin{tcolorbox}[colback=cajafondo,colframe=umagpurple!40,width=0.8\textwidth,arc=3pt,boxrule=0.6pt]
\centering
#2
\end{tcolorbox}

\vspace{1cm}
{\color{umaggray} Versión del documento: \versionmanual}

\vspace{1cm}
\end{titlepage}
}
